\documentclass[landscape]{beamer}
\usetheme{default}
\usepackage{amsmath, amssymb}

\title{CSCE 222 --- Sets (But Actually Understandable)}
\author{You Are Not Bad at Math}
\date{}

\begin{document}

%------------------------------------------------
\begin{frame}
\titlepage
\end{frame}

%------------------------------------------------
\begin{frame}{First: Breathe. This Is Not Hard.}
You are not lost because this is hard.\\
You are lost because it was presented without patterns.

\vspace{0.5cm}

Today we:
\begin{itemize}
\item Strip everything to its core idea
\item Define every symbol like code
\item Connect it to CS thinking
\item Use basketball, games, and logic patterns
\end{itemize}

Sets are just:
\[
\textbf{Organized collections of distinct objects}
\]

That is it.
\end{frame}

%------------------------------------------------
\begin{frame}{CHEAT SHEET (1/4)}
\begin{columns}

\column{0.48\textwidth}
\textbf{Set Definition}

A set is a collection of distinct elements.

Example:
\[
S = \{1,2,3,4,5\}
\]

$x \in S$ means:
\begin{itemize}
\item $x$ = element
\item $\in$ = "is in"
\item $S$ = the set
\end{itemize}

$x \notin S$ means $x$ is not in $S$.

\column{0.48\textwidth}
\textbf{Set Equality}

\[
A = B \iff \forall x (x \in A \leftrightarrow x \in B)
\]

Variables:
\begin{itemize}
\item $A,B$ = sets
\item $x$ = element
\item $\forall$ = for all
\item $\leftrightarrow$ = if and only if
\end{itemize}

Translation:
Every element in $A$ is in $B$ AND vice versa.
Order does not matter.
\end{columns}
\end{frame}

%------------------------------------------------
\begin{frame}{CHEAT SHEET (2/4)}
\begin{columns}

\column{0.48\textwidth}
\textbf{Subset}

\[
A \subseteq B
\]

Definition:
\[
\forall x (x \in A \rightarrow x \in B)
\]

Means:
If $x$ is in $A$, it must be in $B$.

Proper subset:
\[
A \subset B
\]
Same as above AND $A \neq B$.

\column{0.48\textwidth}
\textbf{Empty Set}

\[
\emptyset = \{\}
\]

Cardinality:
\[
|\emptyset| = 0
\]

Be careful:
\[
\{\emptyset\}
\]
is a set with ONE element.

Like:
empty box vs box containing empty box.
\end{columns}
\end{frame}

%------------------------------------------------
\begin{frame}{CHEAT SHEET (3/4)}
\begin{columns}

\column{0.48\textwidth}
\textbf{Operations}

Union:
\[
A \cup B
\]
Elements in A OR B.

Intersection:
\[
A \cap B
\]
Elements in BOTH.

Difference:
\[
B - A
\]
In B but NOT in A.

\column{0.48\textwidth}
\textbf{Complement}

\[
A^c
\]

All elements in universal set $U$
that are NOT in $A$.

Everything depends on $U$.

No universal set = no complement.
\end{columns}
\end{frame}

%------------------------------------------------
\begin{frame}{CHEAT SHEET (4/4)}
\begin{columns}

\column{0.48\textwidth}
\textbf{Power Set}

\[
P(S) = \text{set of all subsets of } S
\]

If:
\[
|S| = n
\]
then:
\[
|P(S)| = 2^n
\]

Pattern:
Each element has 2 choices:
in or out.

\column{0.48\textwidth}
\textbf{Cartesian Product}

\[
A \times B
\]

All ordered pairs $(a,b)$:
\[
a \in A, \quad b \in B
\]

If:
\[
|A|=n, |B|=m
\]

Then:
\[
|A \times B| = nm
\]

Multiplication principle.
\end{columns}
\end{frame}

%------------------------------------------------
\begin{frame}{Set-Builder Notation (Slowly)}
\[
\{x \in S \mid P(x)\}
\]

Break it:

\begin{itemize}
\item $x$ = placeholder variable
\item $S$ = big set we pull from
\item $\mid$ = such that
\item $P(x)$ = condition
\end{itemize}

Example:
\[
\{x \in \mathbb{Z} \mid 2 < x \le 5\}
\]

Step 1: $\mathbb{Z}$ = integers  
Step 2: numbers greater than 2  
Step 3: less or equal to 5  

Result:
\[
\{3,4,5\}
\]

Like filtering a list in code.
\end{frame}

%------------------------------------------------
\begin{frame}{Why Subset Definition Has $\forall x$}
Definition:
\[
\forall x (x \in A \rightarrow x \in B)
\]

Why $\forall$?

Without it, we only check one element.

Why implication $\rightarrow$?

Because we only care about elements in $A$.

If $x \notin A$, rule is irrelevant.

This prevents false failures.

This is defensive programming logic.
\end{frame}

%------------------------------------------------
\begin{frame}{Power Set Pattern (This Is Beautiful)}
Let:
\[
S = \{a,b\}
\]

Subsets:

\[
\emptyset
\]
\[
\{a\}
\]
\[
\{b\}
\]
\[
\{a,b\}
\]

Why $2^2 = 4$?

Each element:
\begin{itemize}
\item include?
\item exclude?
\end{itemize}

Binary thinking:
00, 01, 10, 11

Power set = bitmasking.
\end{frame}

%------------------------------------------------
\begin{frame}{Cartesian Product (Video Game View)}
Let:
\[
A = \{3,5,7\}
\]
\[
B = \{2,4\}
\]

$A \times B$:

Fix 3:
\[
(3,2),(3,4)
\]

Fix 5:
\[
(5,2),(5,4)
\]

Fix 7:
\[
(7,2),(7,4)
\]

This is nested loops.

for a in A:
\quad for b in B:

Total = $3 \cdot 2 = 6$
\end{frame}

%------------------------------------------------
\begin{frame}{Relation = Subset of Product}
Definition:
\[
R \subseteq A \times B
\]

Meaning:
A relation picks SOME pairs.

Example:

If $(x,y) \in R$ means $x<y$

Then we filter product pairs.

Relation = filtered Cartesian product.

Like matchmaking in 2K:
Not everyone plays everyone.
\end{frame}

%------------------------------------------------
\begin{frame}{Worked Example (Activity)}
Universal set:
\[
U=\{1,2,3,4,5,6,7\}
\]

\[
A=\{1,3,5,7\}
\]
\[
B=\{4,5,6,7\}
\]
\end{frame}

%------------------------------------------------
\begin{frame}{Solution Part 1}
Union:
\[
A \cup B
\]

Combine unique elements:
\[
\{1,3,4,5,6,7\}
\]

Intersection:
\[
A \cap B
\]

Common elements:
\[
\{5,7\}
\]
\end{frame}

%------------------------------------------------
\begin{frame}{Solution Part 2}
Difference:
\[
B-A
\]

Elements in B not in A:
\[
\{4,6\}
\]

Complement:
\[
A^c
\]

Everything in $U$ not in $A$:
\[
\{2,4,6\}
\]
\end{frame}

%------------------------------------------------
\begin{frame}{Homework Survival Mindset}
Your professor writes symbols.

You translate:

\begin{itemize}
\item $\cup$ = OR
\item $\cap$ = AND
\item $\subseteq$ = if in A then in B
\item $P(S)$ = binary combinations
\item $\times$ = nested loops
\end{itemize}

Everything reduces to:
\begin{itemize}
\item membership
\item filtering
\item combinations
\end{itemize}

You already do this in CS daily.
\end{frame}

%------------------------------------------------
\begin{frame}{Extra Example (Power Set)}
Let:
\[
S=\{1,2,3\}
\]

How many subsets?

\[
|P(S)| = 2^3 = 8
\]

Why?

Each element:
include or exclude.

Binary:
000
001
010
011
100
101
110
111
\end{frame}

%------------------------------------------------
\begin{frame}{Final Perspective}
Sets are not scary.

They are:

\begin{itemize}
\item Lists without duplicates
\item Logical filters
\item Nested loops
\item Binary choices
\end{itemize}

This entire lecture is:
organized thinking.

If you understand patterns,
you win.

You are absolutely capable of this.
\end{frame}

\end{document}